% answers_short.tex —— 2011 年·全国大纲卷·文科 · 简短答案速查
% 由 tools/make_short_answers.py 自 answers.tex 的 \gksolhead/\gkanswer 提取，
% 并与 2011/2011_数学真题_全国大纲文_答案要点.md 逐题比对后生成，勿手改（改源 MD 后重跑）。

\begin{gknotice}
\textbf{使用说明}：本节只给最终答案便于快速核对，完整推导见解析版。答案为编写者独立推导，\textbf{非官方标准答案}。
\end{gknotice}

\gkpart{选择题}
\begin{center}\begin{tabular}{ccccccccccccc}
\toprule
题号 & 1 & 2 & 3 & 4 & 5 & 6 & 7 & 8 & 9 & 10 & 11 & 12 \\
\midrule
答案 & D & B & B & C & A & D & C & C & B & A & C & D \\
\bottomrule
\end{tabular}\end{center}

\gkpart{填空题}
\begin{itemize}[leftmargin=2.6em]
  \item[\textbf{13.}] $0$
  \item[\textbf{14.}] $-\dfrac1{\sqrt5}$
  \item[\textbf{15.}] $\dfrac23$
  \item[\textbf{16.}] $6$
\end{itemize}

\gkpart{解答题}
\begin{itemize}[leftmargin=2.6em]
  \item[\textbf{17.}] 有两组解：$(a_n,S_n)=\left(2\cdot3^{n-1},\,3^n-1\right)$ 或 $(a_n,S_n)=\left(3\cdot2^{n-1},\,3(2^n-1)\right)$。
  \item[\textbf{18.}] (1) $B=45^\circ$；(2) $a=1+\sqrt3$，$c=\sqrt6$
  \item[\textbf{19.}] (1) 至少购买一种保险的概率为 $0.8$；(2) 恰有 $1$ 位车主两种保险都不购买的概率为 $\binom31(0.2)(0.8)^2=0.384=\dfrac{48}{125}$
  \item[\textbf{20.}] (1) 见解析版；(2) $AB$ 与平面 $SBC$ 所成角的大小满足 $\sin\theta=\dfrac{\sqrt3}{\sqrt7}=\dfrac{\sqrt{21}}7$，即 $\theta=\arcsin\dfrac{\sqrt{21}}7$
  \item[\textbf{21.}] (1) 切线为 $y=(3-6a)x+12a-4$；(2) $a\in\left(-\dfrac52,-1-\sqrt2\right)$
  \item[\textbf{22.}] (1) 见解析版；(2) 见解析版
\end{itemize}
